\documentclass[12pt,a4paper]{article}

\usepackage[a4paper,margin=1in]{geometry}
\usepackage{graphicx}
\usepackage{booktabs}
\usepackage{array}
\usepackage{longtable}
\usepackage{hyperref}
\usepackage{xcolor}
\usepackage{listings}
\usepackage{enumitem}
\usepackage{fancyhdr}
\usepackage{titlesec}

\hypersetup{
    colorlinks=true,
    linkcolor=blue,
    urlcolor=blue,
    citecolor=blue
}

\definecolor{codegray}{RGB}{245,245,245}
\definecolor{codeborder}{RGB}{210,210,210}

\lstdefinestyle{projectcode}{
    backgroundcolor=\color{codegray},
    basicstyle=\ttfamily\small,
    breaklines=true,
    frame=single,
    rulecolor=\color{codeborder},
    columns=fullflexible,
    keepspaces=true,
    showstringspaces=false
}

\pagestyle{fancy}
\fancyhf{}
\lhead{SPE Project Report}
\rhead{Chatbot DevSecOps Pipeline}
\cfoot{\thepage}

\titleformat{\section}{\Large\bfseries}{\thesection}{1em}{}
\titleformat{\subsection}{\large\bfseries}{\thesubsection}{1em}{}

\begin{document}

\begin{titlepage}
    \centering
    \vspace*{1.2cm}
    {\Huge\bfseries Chatbot DevSecOps Pipeline\par}
    \vspace{0.4cm}
    {\Large Software Project Engineering Report\par}
    \vspace{1.2cm}
    {\large Project Repository: \texttt{poojann-pandyaa/chatbot-devsecops-pipeline}\par}
    \vspace{0.8cm}
    {\large Technology Stack: Next.js, Docker, Jenkins, Ansible, Kubernetes, Terraform, Vault, SonarQube, Trivy\par}
    \vfill
    \begin{tabular}{rl}
        Student Name: & \textbf{Poojan Pandya} \\
        Course: & \textbf{SPE} \\
        Submission: & \textbf{Final Project} \\
        Date: & \textbf{\today}
    \end{tabular}
    \vfill
\end{titlepage}

\tableofcontents
\newpage

\section{Abstract}

This project implements a DevSecOps-enabled deployment pipeline for a web-based
AI chatbot. The application is a Next.js frontend and API service that connects
to an OpenAI-compatible large language model provider. The project focuses not
only on application delivery, but also on automation, security scanning,
infrastructure provisioning, secret management, containerization, Kubernetes
deployment, and production-readiness controls.

The final system uses Jenkins for CI/CD orchestration, Docker for packaging,
SonarQube and Trivy for quality and security checks, Ansible for deployment
automation, Terraform for EKS infrastructure, Vault-backed Kubernetes secrets for
API key handling, and Kubernetes manifests for scalable production deployment.
The Kubernetes layer includes namespace isolation, rolling updates, liveness and
readiness probes, Ingress with TLS, a Pod Disruption Budget, and Horizontal Pod
Autoscaling using resource and application-level metrics.

\section{Introduction}

Modern software delivery requires more than writing application code. Reliable
systems must be built, tested, scanned, packaged, deployed, monitored, and
scaled through repeatable automation. This project demonstrates those practices
through a chatbot application deployed using a complete DevSecOps pipeline.

The chatbot application provides a user interface for interacting with a large
language model. The surrounding engineering work transforms the application from
a simple web app into a deployable system with automated quality gates,
container builds, image scanning, Kubernetes rollout verification, and
infrastructure-as-code support.

\section{Project Objectives}

The major objectives of this project are:

\begin{itemize}[leftmargin=*]
    \item Build and package a Next.js chatbot application as a Docker image.
    \item Automate build, scan, image publishing, and deployment through Jenkins.
    \item Integrate static code analysis using SonarQube.
    \item Integrate filesystem and container image vulnerability scanning using Trivy.
    \item Use Ansible as the deployment driver for Kubernetes resources.
    \item Store sensitive API keys through Kubernetes secrets prepared by the Vault configuration role.
    \item Provision cloud Kubernetes infrastructure using Terraform for AWS EKS.
    \item Deploy the application with production-grade Kubernetes controls.
    \item Demonstrate autoscaling using Kubernetes HPA and application metrics.
\end{itemize}

\section{System Overview}

The system contains five main layers:

\begin{enumerate}[leftmargin=*]
    \item \textbf{Application Layer:} A Next.js chatbot UI and API service.
    \item \textbf{Container Layer:} A multi-stage Dockerfile builds and runs the app.
    \item \textbf{CI/CD Layer:} Jenkins automates dependency installation, analysis, scanning, build, push, and deployment.
    \item \textbf{Deployment Automation Layer:} Ansible applies Kubernetes manifests and verifies rollout status.
    \item \textbf{Infrastructure Layer:} Terraform defines AWS EKS cluster resources.
\end{enumerate}

\subsection{High-Level Architecture}

\begin{center}
\begin{tabular}{>{\raggedright\arraybackslash}p{4cm} >{\raggedright\arraybackslash}p{9cm}}
\toprule
\textbf{Component} & \textbf{Responsibility} \\
\midrule
Next.js App & Provides chatbot UI, chat API, health endpoint, and metrics endpoint. \\
Docker & Packages the application into a portable container image. \\
Jenkins & Runs CI/CD stages including build, security checks, image push, and Ansible deployment. \\
SonarQube & Performs code quality and maintainability analysis. \\
Trivy & Scans source files and Docker images for vulnerabilities. \\
Ansible & Applies Kubernetes resources and waits for deployment rollout completion. \\
Kubernetes & Runs the chatbot workload with replicas, probes, service, Ingress, PDB, and HPA. \\
Vault Role & Prepares secret flow and creates the Kubernetes secret used by the chatbot pod. \\
Terraform & Defines AWS EKS cluster, node group, IAM, VPC, and provider configuration. \\
\bottomrule
\end{tabular}
\end{center}

\section{Application Design}

The application is implemented using Next.js. It contains a browser-based chat
interface and API routes under \texttt{app/pages/api}. The main chat route
streams responses from an OpenAI-compatible API endpoint. In this deployment,
the Kubernetes manifest configures the host as:

\begin{lstlisting}[style=projectcode]
OPENAI_API_HOST=https://api.groq.com/openai
OPENAI_API_MODEL=llama-3.3-70b-versatile
\end{lstlisting}

The API key is provided through the environment variable
\texttt{OPENAI\_API\_KEY}. In Kubernetes, this value is injected from the
\texttt{chatbot-secrets} secret.

\subsection{Health and Metrics Endpoints}

Two operational endpoints were added to support production deployment:

\begin{itemize}[leftmargin=*]
    \item \texttt{/api/health}: returns a simple JSON status response for Kubernetes liveness and readiness probes.
    \item \texttt{/api/metrics}: exposes Prometheus-compatible metrics, including \texttt{http\_requests\_per\_second}, \texttt{http\_active\_requests}, and \texttt{http\_requests\_total}.
\end{itemize}

These endpoints make the application observable and allow Kubernetes to
distinguish between a running container and a healthy application.

\section{Dockerization}

The application uses a multi-stage Dockerfile. The stages are:

\begin{enumerate}[leftmargin=*]
    \item \textbf{Base:} sets up the Node.js working directory and copies package manifests.
    \item \textbf{Dependencies:} runs \texttt{npm ci} for deterministic dependency installation.
    \item \textbf{Build:} copies the application source and runs \texttt{npm run build}.
    \item \textbf{Production:} copies the compiled Next.js application and starts it with \texttt{npm start}.
\end{enumerate}

This keeps the runtime image focused on production execution and ensures that
build steps are reproducible.

\section{CI/CD Pipeline}

The Jenkins pipeline is defined in the root \texttt{Jenkinsfile}. It performs
the following stages:

\begin{longtable}{>{\raggedright\arraybackslash}p{4cm} >{\raggedright\arraybackslash}p{9cm}}
\toprule
\textbf{Stage} & \textbf{Purpose} \\
\midrule
Checkout SCM & Pulls the source code and calculates a short Git commit image tag. \\
Install Dependencies & Installs application dependencies inside the \texttt{app} directory. \\
SonarQube Analysis & Runs static analysis through SonarQube. \\
Quality Gate & Waits for the SonarQube quality gate result. \\
OWASP FS Scan & Placeholder stage for OWASP dependency scanning once the NVD key is available. \\
Trivy FS Scan & Scans the source tree and dependencies for vulnerabilities. \\
Docker Build and Push & Builds the Docker image and pushes both \texttt{:latest} and commit-tagged images. \\
Trivy Image Scan & Scans the commit-tagged Docker image. \\
Deploy to Container & Runs the commit-tagged image as a local container for quick validation. \\
Ansible Deploy & Runs the Ansible playbook to apply Kubernetes resources and verify rollout status. \\
\bottomrule
\end{longtable}

\subsection{Traceable Image Tags}

The pipeline generates a short image tag from the current Git commit:

\begin{lstlisting}[style=projectcode]
env.IMAGE_TAG = sh(script: 'git rev-parse --short=7 HEAD',
                   returnStdout: true).trim()
\end{lstlisting}

The Docker image is pushed using both the immutable commit tag and the latest
tag:

\begin{lstlisting}[style=projectcode]
docker buildx build \
  --platform linux/arm64 \
  -t ${IMAGE_NAME}:${IMAGE_TAG} \
  -t ${IMAGE_NAME}:latest \
  --push app/
\end{lstlisting}

This improves traceability because each Kubernetes rollout can be mapped back to
a specific source code commit.

\section{Security Engineering}

Security is integrated at multiple stages of the project.

\subsection{Static Analysis}

SonarQube is used in the Jenkins pipeline to analyze code quality, maintainable
design, bugs, and code smells. The pipeline includes a quality gate stage, which
allows the build to react to the quality result.

\subsection{Vulnerability Scanning}

Trivy is used twice:

\begin{itemize}[leftmargin=*]
    \item \textbf{Filesystem scan:} \texttt{trivy fs . > trivyfs.txt}
    \item \textbf{Image scan:} \texttt{trivy image \$\{IMAGE\_NAME\}:\$\{IMAGE\_TAG\} > trivyimage.txt}
\end{itemize}

The filesystem scan checks project files and dependency manifests. The image
scan checks the final deployable container image. This ensures that both source
dependencies and runtime image packages are inspected.

\subsection{Secret Management}

The chatbot requires an API key to call the LLM provider. The key is not
hard-coded into the application or Docker image. Instead, Kubernetes injects it
from the secret named \texttt{chatbot-secrets}.

Example command:

\begin{lstlisting}[style=projectcode]
kubectl create secret generic chatbot-secrets \
  --from-literal=OPENAI_API_KEY='gsk_your_key_here' \
  -n chatbot-prod \
  --dry-run=client -o yaml | kubectl apply -f -
\end{lstlisting}

The Ansible Vault configuration role also creates the Kubernetes secret in the
target namespace using the provided \texttt{openai\_api\_key} variable.

\section{Infrastructure as Code}

The \texttt{EKS-TF} directory contains Terraform files for AWS EKS
infrastructure. The infrastructure code includes:

\begin{itemize}[leftmargin=*]
    \item AWS provider configuration.
    \item EKS cluster definition.
    \item EKS node group definition.
    \item IAM roles and policies.
    \item VPC and subnet-related configuration.
    \item Terraform backend and variables.
\end{itemize}

Using Terraform makes the infrastructure reproducible. Instead of manually
creating Kubernetes clusters, the EKS environment can be declared, reviewed,
versioned, and recreated.

\section{Ansible Deployment Automation}

Ansible is used as the actual Kubernetes deployment driver. The playbook
\texttt{ansible/site.yml} includes three roles:

\begin{itemize}[leftmargin=*]
    \item \texttt{vault-config}: prepares Vault-related configuration and Kubernetes secrets.
    \item \texttt{k8s-deploy}: applies Kubernetes manifests and verifies rollout status.
    \item \texttt{docker}: legacy role for direct Docker deployment.
\end{itemize}

The Jenkins deployment stage calls:

\begin{lstlisting}[style=projectcode]
ansible-playbook -i ansible/inventory/hosts.ini ansible/site.yml \
  --extra-vars "image_tag=${IMAGE_TAG} git_commit=${GIT_COMMIT} namespace=${K8S_NAMESPACE}"
\end{lstlisting}

This design avoids using Jenkins as a raw shell wrapper around \texttt{kubectl}.
Instead, Jenkins delegates deployment logic to Ansible, which makes the
deployment repeatable and easier to extend.

\section{Kubernetes Deployment Design}

The Kubernetes deployment was upgraded from a basic manifest into a more
production-ready setup. The resources are stored under the \texttt{k8s}
directory.

\subsection{Namespace Isolation}

The application runs in the dedicated namespace \texttt{chatbot-prod}. This
separates chatbot resources from default cluster resources and supports cleaner
access control, observability, and deletion.

\subsection{Deployment Strategy}

The deployment uses two replicas and a rolling update strategy:

\begin{lstlisting}[style=projectcode]
replicas: 2
strategy:
  type: RollingUpdate
  rollingUpdate:
    maxSurge: 1
    maxUnavailable: 0
\end{lstlisting}

The value \texttt{maxUnavailable: 0} prevents Kubernetes from intentionally
taking an old pod down before a replacement is ready. This supports
zero-downtime deployment.

\subsection{Health Probes}

The pod defines both liveness and readiness probes against
\texttt{/api/health}. The readiness probe controls whether a pod receives
traffic. The liveness probe allows Kubernetes to restart the container if the
application becomes unhealthy.

\subsection{Service and Ingress}

The service is a \texttt{ClusterIP} service named \texttt{chatbot-service}.
External traffic is intended to enter through Nginx Ingress:

\begin{itemize}[leftmargin=*]
    \item Host: \texttt{chatbot.local}
    \item TLS secret: \texttt{chatbot-tls}
    \item Backend service: \texttt{chatbot-service}
    \item Backend port: \texttt{80}
\end{itemize}

This is closer to real production Kubernetes than directly exposing every
application through a \texttt{LoadBalancer} service.

\subsection{Pod Disruption Budget}

The Pod Disruption Budget requires at least one chatbot pod to remain available:

\begin{lstlisting}[style=projectcode]
apiVersion: policy/v1
kind: PodDisruptionBudget
metadata:
  name: chatbot-pdb
spec:
  minAvailable: 1
\end{lstlisting}

This protects the application during voluntary disruptions such as node
maintenance.

\section{Autoscaling and Observability}

The Horizontal Pod Autoscaler is configured with CPU, memory, and pod-level
request-rate metrics:

\begin{itemize}[leftmargin=*]
    \item CPU target utilization: 50 percent.
    \item Memory target utilization: 80 percent.
    \item Pod metric: \texttt{http\_requests\_per\_second}, average value 10.
\end{itemize}

The application exposes \texttt{/api/metrics}, which produces
Prometheus-compatible metric output. The pod template includes scrape
annotations:

\begin{lstlisting}[style=projectcode]
prometheus.io/scrape: "true"
prometheus.io/path: "/api/metrics"
prometheus.io/port: "3000"
\end{lstlisting}

For CPU and memory scaling, the Kubernetes Metrics Server is required. For the
custom request-rate metric, Prometheus and Prometheus Adapter are required so
that Kubernetes HPA can consume the application metric.

\subsection{Scaling Demonstration Commands}

The following commands can be used to observe scaling:

\begin{lstlisting}[style=projectcode]
kubectl get hpa -n chatbot-prod -w
kubectl get pods -n chatbot-prod -w
kubectl describe hpa chatbot-hpa -n chatbot-prod
\end{lstlisting}

A simple load generator can be started inside the cluster:

\begin{lstlisting}[style=projectcode]
kubectl run load-generator \
  --rm -i --tty \
  --image=busybox \
  -n chatbot-prod \
  -- /bin/sh
\end{lstlisting}

Inside the shell:

\begin{lstlisting}[style=projectcode]
while true; do wget -q -O- http://chatbot-service/api/health; done
\end{lstlisting}

\section{Testing and Verification}

The Docker image was built successfully using the project Dockerfile. The build
completed the Next.js production compilation and showed the API routes,
including:

\begin{itemize}[leftmargin=*]
    \item \texttt{/api/chat}
    \item \texttt{/api/health}
    \item \texttt{/api/metrics}
    \item \texttt{/api/models}
\end{itemize}

The application container was started locally with Docker on port 3000. The
health and metrics endpoints were verified from inside the running container:

\begin{lstlisting}[style=projectcode]
/api/health  -> {"status":"ok"}
/api/metrics -> Prometheus text metrics
\end{lstlisting}

YAML syntax was also checked for the Kubernetes and Ansible YAML files using a
local parser. A full \texttt{ansible-playbook --syntax-check} could not be run
in the local shell because \texttt{ansible-playbook} was not installed in the
environment.

\section{Risk Analysis}

\begin{longtable}{>{\raggedright\arraybackslash}p{4cm} >{\raggedright\arraybackslash}p{5cm} >{\raggedright\arraybackslash}p{4cm}}
\toprule
\textbf{Risk} & \textbf{Impact} & \textbf{Mitigation} \\
\midrule
Invalid API key & Chatbot cannot call the LLM provider. & Use Kubernetes secret and verify key before rollout. \\
Vulnerable dependencies & Security exposure in app or image. & Run Trivy filesystem and image scans in Jenkins. \\
Failed deployment & Users may receive errors after deployment. & Use readiness probes and \texttt{kubectl rollout status}. \\
Pod downtime during update & Temporary service outage. & Use two replicas and \texttt{maxUnavailable: 0}. \\
Node maintenance disruption & Reduced availability. & Use Pod Disruption Budget with \texttt{minAvailable: 1}. \\
Unconfigured custom metrics & HPA may not scale on request rate. & Install Prometheus and Prometheus Adapter before relying on custom metric. \\
\bottomrule
\end{longtable}

\section{Project Outcomes}

The project demonstrates a full DevSecOps workflow:

\begin{itemize}[leftmargin=*]
    \item Application code is built and packaged into Docker images.
    \item Images are tagged with Git commit hashes for traceability.
    \item Static analysis and vulnerability scanning are part of the pipeline.
    \item Secrets are separated from source code and injected at runtime.
    \item Kubernetes deployment uses production controls rather than a minimal demo manifest.
    \item Ansible performs deployment automation and rollout verification.
    \item Terraform provides infrastructure-as-code support for AWS EKS.
    \item Health and metrics endpoints support reliability and observability.
\end{itemize}

\section{Future Enhancements}

The project can be improved further with:

\begin{itemize}[leftmargin=*]
    \item Prometheus Adapter configuration committed as code for custom HPA metrics.
    \item Grafana dashboards for chatbot traffic, latency, errors, and pod resource usage.
    \item External Secrets Operator or Vault Agent Injector for deeper Vault integration.
    \item Argo CD or Flux for GitOps-based continuous deployment.
    \item Jenkins quality gates configured to fail builds on critical vulnerabilities.
    \item Automated integration tests against the deployed Kubernetes service.
    \item Production TLS certificate automation through cert-manager.
\end{itemize}

\section{Conclusion}

This SPE project successfully converts a chatbot web application into a
DevSecOps-ready system. The result includes CI/CD automation, image scanning,
code quality analysis, containerization, Kubernetes deployment, secret handling,
infrastructure-as-code, health checks, metrics, and autoscaling configuration.

The most important engineering improvement is the shift from a basic deployment
to a traceable and verifiable release process. Jenkins builds a commit-tagged
image, Ansible applies Kubernetes resources, and Kubernetes verifies rollout
health using probes and rollout status. These practices make the project closer
to real production software delivery and demonstrate strong understanding of
modern DevOps and DevSecOps principles.

\end{document}
